% cap_07_conclusao.tex  –  Capítulo 7
% Conclusão
\chapter{Conclusão}
\label{ch:conclusao}

Esta tese introduziu as Redes de Petri Hierárquicas de Sinais, uma extensão
formal da teoria clássica de Redes de Petri que unifica modelagem metabólica e
regulatória por meio de semântica de consumo de sinal e organização hierárquica
em camadas.

\section{Contribuições}
\label{sec:contribuicoes}

\begin{sloppypar}
A primeira contribuição é o \emph{formalismo} apresentado no
Capítulo~\ref{ch:shpn}. A 13-tupla das Redes de Petri Hie\-rár\-qui\-cas
de Sinais
estende a 5-tupla clássica acrescentando lugares de sinal
($\Psi \subset P$), arcos de fluxo de sinal ($F_s$) e camadas hierárquicas
($\lambda$), sem remover nenhum elemento da definição original e preservando,
portanto, compatibilidade retroativa: quando $\Psi = \emptyset$, recupera-se a
semântica clássica integralmente. O parâmetro termodinâmico $\Gamma = (K, n,
\varepsilon)$ fundamenta os limiares de habilitação em constantes cinéticas
mensuráveis, removendo o ajuste arbitrário e ancorando a decisão regulatória em
quantidades acessíveis a determinação experimental independente.

A segunda contribuição é o conjunto de \emph{propriedades teóricas demonstradas}
no mesmo capítulo. A aciclicidade da hierarquia de sinais, a garantia de
preempção, a formação de fronteiras de bacia e a independência fraca constituem
as garantias formais que viabilizam a predição de limiares por análise de
alcança\-bilidade. A combinação dessas propriedades reduz a complexidade da
análise de exponencial a polinomial em modelos com hierarquia
de\-com\-po\-ní\-vel, abrindo caminho para tratamento de modelos de escala
ge\-nô\-mi\-ca que escapariam das técnicas clássicas.

A terceira contribuição é a \emph{validação biológica} apresentada no
Capítulo~\ref{ch:validacao-bacillus}. O modelo de esporulação do
\textit{Bacillus subtilis} (com 40 lugares, incluindo
9 lugares de parâmetro $\square$ e 3 lugares de sinal espacial, e 36
transições) foi parametrizado exclusivamente a partir de dados
cinéticos publicados e produziu a separatriz de comprometimento emergente de
$\sigma^H$ a $\theta_{\sigma H} \approx 1{,}60\,\mu$M sob varredura fatorial
$2 \times 3 \times 4$ em disponibilidade de substrato, dose de Spo0A* e SinR
(25~condições, 1250~simulações estocásticas). A separatriz reproduz a
assimetria de \textcite{Fujita2005} --- acumulação gradual via fosforretransmissão
(${\approx}\,48\%$ de réplicas esporulantes em inanição) versus indução constitutiva
abrupta de Spo0A* em meio rico ($0$--$6\%$) --- sem ajuste de nenhum parâmetro
ao resultado experimental.
A robustez da separatriz através das 25~condições confirma a
preditividade do formalismo.

A quarta contribuição é a \emph{plataforma computacional} descrita no
Capítulo~\ref{ch:implementacao}. A arquitetura multi-painel, o duplo motor de
simulação --- equações diferenciais ordinárias resolvidas por Runge-Kutta de
quarta ordem acopladas a $\tau$-leaping estocástico por divisão de operador
--- e o escalonamento hierárquico, somados à importação automatizada a partir
de KEGG e BRENDA, traduzem o formalismo em ferramenta executável de código
aberto. A plataforma é, em si, um veículo de validação: cada capacidade
exposta na tese corresponde a um caminho efetivamente percorrido durante a
construção dos modelos.
\end{sloppypar}

\section{Limitações}
\label{sec:limitacoes-extensoes}

A validação aqui apresentada não esgota o programa de pesquisa. A
\emph{disponibilidade de parâmetros} permanece como restrição prática: a
automação contra a base BRENDA recupera cerca de oitenta e cinco por cento
das constantes cinéticas para vias bem caracterizadas, mas redes pouco
estudadas continuam a exigir parametrização experimental específica. A
análise de sensibilidade demonstrou que as predições toleram a incerteza
realista observada nessas bases, porém a abordagem ainda requer conhecimento
das taxas com precisão de pelo menos uma ordem de magnitude.

O \emph{escopo de validação} concentrou-se na esporulação bacteriana, um
sistema unicelular procariótico, justamente para alcançar precisão
quantitativa em um caso bem caracterizado. Sistemas eucarióticos, com
remodelação de cromatina e regulação por RNAs longos não codificantes,
desenvolvimento multicelular, com gradientes de morfogênio e sinalização
do tipo Notch, e mecanismos não metabólicos como a mecanotransdução
permanecem como territórios a serem percorridos, e cada um trará desafios
formais específicos que poderão refinar a teoria.

A \emph{complexidade computacional} estabelece um limite teórico inevitável:
a análise de alcançabilidade é Ackermann-completa~\cite{Czerwinski2021} no pior caso. A
independência fraca e a decomposição por camada mitigam esse custo na
maioria dos modelos biologicamente plausíveis, mas hubs regulatórios em
redes realistas criam gargalos de análise que recomendam o desenvolvimento
de métodos aproximados, equilibrando precisão e tratabilidade.

Finalmente, a plataforma carece de \emph{verificação formal} integrada por
lógica temporal computacional ou por lógica temporal linear. A integração
com ferramentas especializadas de verificação de Redes de Petri
complementaria a validação baseada em simulação ao oferecer garantias
exaustivas sobre propriedades temporais de comprometimento e
irreversibilidade.

\section{Direções Futuras}
\label{sec:direcoes-futuras}

A continuação natural deste trabalho aponta para várias frentes
complementares. A \emph{extensão multicelular} introduzirá lugares de sinal
localizados com dinâmicas de difusão, permitindo modelar gradientes de
morfogênio, sinalização Notch entre células vizinhas e mecanismos de
sensoriamento de quórum em populações bacterianas. A camada hierárquica $\lambda$,
hoje atribuída manualmente, pode beneficiar-se de \emph{aprendizado
de máquina}: redes neurais sobre grafos treinadas em bases de vias
anotadas automatizariam a atribuição da camada e libertariam o pesquisador
da escolha topológica, hoje guiada por intuição biológica.

Em outra direção, há o \emph{embasamento termodinâmico} mais profundo do
formalismo: conectar a semântica de consumo de sinal à dissipação de energia
livre fundamentaria a regulação na física estatística. Conjectura-se que a
forma funcional de $\Gamma$ possa, em condições a serem estabelecidas, ser
derivada de princípios de produção mínima de entropia em vez de tomada como
postulado; essa investigação permanece como programa para trabalho futuro. A \emph{validação ampla} já
contempla candidatos imediatos --- persistência bacteriana sob estresse
antibiótico, quiescência de leveduras em fase estacionária, formação de
\textit{dauer} em \textit{Caenorhabditis elegans} e pontos de checagem do
ciclo celular em mamíferos --- todos sistemas em que decisões dependentes
de energia possuem limiares quantificáveis na literatura experimental.

A \emph{interoperabilidade plena com SBML} é o complemento natural da
capacidade de importação descrita no Capítulo~\ref{ch:implementacao}: hoje
a SHYpn consome modelos curados do BioModels mapeando-os para a 13-tupla
SHPN acrescida do grafo de parâmetros, mas o caminho inverso --- exportar
um modelo SHPN como documento SBML --- ainda não está implementado. A
exportação inevitavelmente perde anotações que SBML não codifica
nativamente ($\Psi$, $F_s$ vs.\ $F_t$, $\lambda$, $\Gamma$, distinção
rede-objeto vs.\ plano experimental); contudo, fornecer um exportador que
preserve essas informações como anotações SBO/MIRIAM padronizadas
permitiria que modelos validados em SHYpn fossem depositados em
BioModels e reutilizados por outros simuladores, fechando o ciclo de
troca entre laboratórios.

Finalmente, a \emph{verificação por lógica temporal} demanda integração da
plataforma com verificadores de modelos estabelecidos, permitindo certificar
exaustivamente propriedades como ``toda trajetória que cruza
$\theta_{\text{eff}}$ atinge eventualmente o estado comprometido'' --- uma
garantia que a simulação por amostragem aproxima mas não estabelece em
sentido formal.

\section{Perspectiva Final}
\label{sec:perspectiva-final}

A biologia de sistemas busca compreensão quantitativa preditiva por meio
de estruturas formais capazes de capturar princípios organizadores
biológicos. As Redes de Petri Hierárquicas de Sinais contribuem para esse
empreendimento formalizando dois conceitos pervasivos porém anteriormente
informais --- o consumo de sinal e a organização hierárquica em camadas --- e
demonstrando que modelos formais podem alcançar predição genuína a partir da
topologia da rede e de parâmetros independentes do fenômeno predito. A
jornada do conceito teórico à plataforma computacional validada ilustra
como métodos formais avançam de abstração matemática para ferramenta
preditiva quando princípios organizadores biológicos informam as estruturas
matemáticas, e não o contrário. Desafios significativos permanecem, mas a
capacidade demonstrada de predizer limiares de decisão quantitativamente
sugere que a fundação teórica aqui estabelecida merece investigação
aprofundada.
sustentada.